\section*{Non-negative rank (author source p. 13:2)}
Given an $m\times n$ matrix $M$ with real non-negative entries, its non-negative rank, $\rk_+(M)$, is the smallest $s$ such that $M$ can be written as
\[M=LR,\]
where $L$ and $R$ are non-negative matrices of dimensions $m\times s$ and $s\times n$, respectively.
\paragraph{Notation.} All logarithms are in base 2 and $[n]:=\{1,\ldots,n\}$.
\section*{6.2 Submatrices of large non-negative rank (author source pp. 13:8--13:9)}
In order to apply Theorem 13, we also need a connection between extension complexity and non-negative rank. This is provided by the notion of slack matrix introduced in [20]. Following [20,2], we now define what it is. Let $V$ be a sequence $v_1,\ldots,v_{m_1}$ of points in $\RR^r$ and $L(x)$ a system $\ell_1(x)\ge b_1,\ldots,\ell_{m_2}(x)\ge b_{m_2}$ of inequalities in $\RR^r$. The slack matrix with respect to $V$ and $L(x)$ is the $m_2\times m_1$ matrix $S$ such that
\[S_{i,j}=\ell_i(v_j)-b_i.\]
Let $P_0:=\conv(V)$ be the convex hull of $V$ and $P_1:=\{x\in\RR^r:L(x)\text{ holds}\}$. If $P_0\subseteq P_1$ then $S$ is non-negative. In [2], we can find:
\begin{originalresult}{Lemma 14 ([2])}
Let $P_0\subseteq P_1$ and $S$ be as above. Define $\xc(P_0,P_1)$ as the minimum $\xc(P)$ over all polyhedra with $P_0\subseteq P\subseteq P_1$. Then
\[\rk_+S-1\le\xc(P_0,P_1)\le\rk_+S.\]
\end{originalresult}
